\documentclass{article}

\textwidth=6.5in
\textheight=8.5in
\voffset=-54pt
\oddsidemargin=0in
\evensidemargin=0in
\setlength{\parskip}{8pt}
\setlength{\parindent}{0pt}

\usepackage{workbook}

\usepackage[workshop]{handout}
\renewcommand{\workshopnote}{CA Notes.} 
\usepackage{lipsum}
 \usepackage{microtype}
 \usepackage{vwcol}

\begin{document}
\htitle{Workshop 11: Slicing to Solve Density Problems, Intro to Differential Equations}
\begin{workshop}
    Welcome to the last workshop of the semester: Workshop 11! The goals for this workshop are for students to:
    \begin{itemize}
    \item use slicing and some geometry to solve density problems with circles and cones.
    \item review the highlights of our introduction to differential equations:
    \begin{itemize}
        \item A differential equation is an equation that involves an unknown function and one or more of its derivatives.
        \item A solution to a differential equation is a function that makes the differential equation true.
        \item In words, the statement $\frac{dP}{dt}=ky$ means that the rate of change of $P$ is proportional to the function $P$ itself by some proportionality constant $k$.
        \item For our purposes, we can check if a function solves a differential equation by (1) taking the derivative of the function and (2) plugging in the original function into the provided differential equation, and seeing if these two quantities match.
    \end{itemize}
    \end{itemize}
\end{workshop}
\begin{enumerate}

\item
 
  City Hall is located in the center of a city. The Financial District is the region within a $3$ km radius of City Hall, and the Residential District is the
  region between $3$ and $9$ km from City Hall.  The population density in
  the Residential District varies with $x$, the distance (in km) from the
  Financial District, and is given by $\rho(x) = 150 e^{ 0.01x}$ people
  per square kilometer.  We'd like to find the number of people living in
  the Residential District.

  \begin {enumerate}

  \item
    \begin{workshop}
      They may find it hard to understand what $x$ means here, since we've
      never done a problem quite like this.  You may need to clarify that
      ``distance to the Financial District'' means ``distance to the
      closest point in the Financial District'', so all points on the
      border of the Financial District have $x = 0$. 
    \end{workshop}

    \begin{problem}
       Make a sketch of the situation
      that shows the meaning of $x$.  Where is the population the most
      dense?  How dense is it there?  Where is the population the least
      dense?  How dense is it there? 
    \end{problem}

    \begin{solution}
      Here's a sketch (the dot in the middle is City Hall, and the gray
      area is the Residential District; the white area in the middle is the
      Financial District).  
      \begin{center}
        \begin{tikzpicture}[rotate=30]
          \begin{axis}[xmin=0, xmax=10, ymin=0, ymax=0, clip=false, x=0.2cm,
            y=0.2cm, hide y axis, xtick={3.2, 9.1},
            xticklabels={0, 6}]
            \draw[fill=lightgray, fill opacity=0.5] (0, 0) circle (9);
            \draw[fill=white] (0, 0) circle (3);
            \filldraw[black] (0, 0) circle(1.5pt);
            
          \end{axis}
        \end{tikzpicture}
      \end{center}
      The population is most dense at the outer edge of the Residential
      District; the density there is $\rho(6) = 150 e^{0.06}$ people per
      square km.  The population is least dense where the Financial
      District meets the Residential District; the density there is $\rho(0) =
      150$ people per square km.
    \end{solution}
\pspace{\vfill}

  \item
    \begin{workshop}
      The notational convention we use is that we always call the number 
      of slices $n$.  When they write about the $k$-th slice, they should
      be writing in terms of $x_k$ and $\Delta x$.  One thing they will
      probably be especially uncomfortable with is the fact that the
      radius of the $k$-th slice is $x_k + 3$, not $x_k$.

      If students wonder about how to find the area of the $k$-th slice,
      the idea we taught them in class is to imagine unrolling the ring to
      get something close to a rectangle with thickness $\Delta x$ and
      width equal to the circumference of the ring.
    \end{workshop}

    \begin{problem}
      In your sketch above, slice the Residential District such that the slices have approximately constant population density. Approximate the number
      of people living in the $k$-th slice.
    \end{problem}

    \begin{solution}

            \begin{center}
          \begin{tikzpicture}
            \begin{axis}[x=0.007in, y=0.007in, xmin=0, xmax=115,
              ymin=-100, ymax=100, axis on top=true, clip=false, hide y
              axis, xtick={37.6, 75, 100}, xticklabels={$x_0$, $x_k$, $x_n$},
              xlabel=$x$, tick label style={fill=none}]      
              \draw[fill=yellow, draw=none] (0, 0) circle(87.5); %
              \draw[fill=white, draw=none] (0, 0) circle (75); %
              \pgfplotsforeachungrouped \x in {37.5, 50, 62.5, 75, ..., 100.1}
              {
                \edef\temp{
                  \noexpand\draw (0, 0) circle(\x);
                }
                \temp
              }         
            \end{axis}
          \end{tikzpicture}
        \end{center}
      Slice the Residential District into concentric circles; the number of people in the $k$-th slice is
      $\approx \rho(x_k) \cdot 2 \pi (x_k + 3) \Delta x$.
    \end{solution}
\pspace{\vfill}
  \item
    \begin{workshop}
      This is always $\displaystyle \sum_{k=1}^n$ of whatever their
      approximation for the $k$-th slice was.
    \end{workshop}

    \begin{problem}
      Write a Riemann sum that approximates the number of people
      living in the Residential District.
    \end{problem}

    \begin{solution}
      $\displaystyle \sum_{k=1}^n \rho(x_k) \cdot 2 \pi (x_k + 3) \Delta
      x$. 
    \end{solution}
    \pspace{\vfill}
    \pspace{\newpage}
    
  \item
    \begin{problem}
      Write but do not evaluate a definite integral that gives the exact number of people living in the
      Residential District.
    \end{problem}

    \begin{solution}
      $\displaystyle \int_0^6 \rho(x) \cdot 2 \pi (x + 3)\,dx$. 
    \end{solution}
\pspace{\vfill}
  \item

    \begin{workshop}
      If we split the Residential Zone into two rings, one with the people
      living within 6 km of City Hall and one with the people living more
      than 6 km away, the inner ring has both a lower population density
      and a smaller area.
    \end{workshop}

    \begin{problem}
      Consider the number of people who live within 6 km of City Hall.  Is
      this less than, greater than, or equal to half the number of people
      living in the Residential District?  How can you answer this question
      without any computation?
    \end{problem}

    \begin{solution}
      \fbox{Less than.} Fewer people live in between $3$ km and $6$ km of City Hall than in between the $6$ km and $9$ km of City Hall. Both the area and the population density between $6$ km and $9$ km of City Hall are greater than the area and the population density between $3$ km and $6$ km of City Hall.
    \end{solution}
\pspace{\vfill}
  \item
    \begin{problem}
      Write but do not evaluate a definite integral that gives the number of people who live within 6
      km of City Hall.  
    \end{problem}

    \begin{solution}
      $\displaystyle \int_0^3 \rho(x) \cdot 2 \pi (x + 3)\,dx$. 
    \end{solution}

  \end{enumerate}
  \pspace{\vfill}
\pspace{\newpage}
\item 

  \begin{workshop}
    They should have done a problem exactly like this in class; the only
    reason it's here is to reinforce the idea that you don't always slice
    a circular region into rings.
  \end{workshop}

  \begin{problem}
    A circular pond has a radius of 20 yards.  A bridge runs
    along a diameter of the pond where tourists stand and feed the fish.  Therefore, the density of fish in
    the pond depends on the distance from the bridge.  The
    density of fish in the pond at a point $x$ yards from the bridge is
    given by $\rho(x) = 40 - x$ fish per square yard.  Write but do not evaluate a definite integral that gives the number of
    fish in the pond.
  \end{problem}
  

\begin{solution}
       
      As usual, we want to slice the region so that the density in each
      slice is approximately constant.  In this case, density varies with
      distance to the bridge, so we'll need to slice parallel to the bridge.
      Let's visualize the bridge running vertically through the center of
      the pond:
      \begin{center}
        \begin{tikzpicture}
          \draw (0, 0) circle (1);
          \draw[blue, thick] (0, -1) -- (0, 1);
        \end{tikzpicture}
      \end{center}
      By symmetry, the left and right halves of the pond have the same
      amount of fish, so we can just focus on the right half of the pond.
      In fact, symmetry tells us that, if we cut the pond both
      horizontally and vertically through the center, each quarter has the
      same amount of fish, so let's just focus on the top right quadrant
      of the pond.  Remembering that $x$ measures the distance from the
      bridge, here is a picture of the quarter pond with the $x$-axis:
      \begin{center}
        \begin{tikzpicture}
          \begin{axis}[xmin=0, xmax=110, ymin=0, ymax=100, x=0.012in,
            y=0.012in, hide y axis, axis x line=bottom, clip=false,
            xlabel=$x$, xtick={0, 100}, xticklabels={0, 20}]  
            \begin{scope}[yshift=8pt]
              \draw (0, 0) -- (100, 0) arc[radius=100, start angle=0, end
              angle=90]; %
              \draw[blue, thick] (0, 0) -- (0, 100);
            \end{scope}
          \end{axis}
        \end{tikzpicture}
      \end{center}
      Slicing parallel to the bridge means slicing the interval $[0, 20]$;
      as usual, we'll slice into $n$ pieces of equal thickness $\Delta x$
      and label the values where we sliced as $x_0 = 0, x_1, \ldots, x_n =
      20$.  Let's focus on the $k$-th slice:
      \begin{center}
        \begin{tikzpicture}
          \begin{axis}[xmin=0, xmax=22, ymin=0, ymax=20, x=0.06in,
            y=0.06in, hide y axis, axis x line=bottom, clip=false,
            xlabel=$x$, xtick={0, 2.5, 5, ..., 20.1},
            xticklabels={$x_0$,,,,$x_k$,,,,$x_n$}] 
            \begin{scope}[yshift=8pt]
            \addplot+[domain=7.5:10, fill=cyan, draw=none,yshift=-8pt]
              {sqrt(20^2-x^2)} -- (10, 0) -- (7.5, 0);
              \draw (0, 0) -- (20, 0) arc[radius=20, start angle=0, end
              angle=90]; %
              \draw[blue, thick] (0, 0) -- (0, 20);
              \pgfplotsforeachungrouped \x in {2.5, 5, ..., 19.9}
              {
                \edef\temp{
                  \noexpand\draw (\x, 0) -- (\x, {sqrt(20^2-\x^2)});
                }
                \temp
              }              
            \end{scope}
          \end{axis}
        \end{tikzpicture}
      \end{center}
      The number of fish in the $k$-th slice is approximately $(\text{fish density
        in $k$-th slice})(\text{area of $k$-th slice})$.  The fish
      density in the $k$-th slice is approximately $\rho(x_k)$.  To
      approximate the area of the $k$-th slice, we can approximate the
      slice as a rectangle of width $\Delta x$.  However, we still need to
      express its height in terms of $x_k$.  The key is to use the shape
      of the pond: it's part of a circle, so we need to use the radius.
      If we draw in the radius (shown in red) below, then we can express
      the height $h_k$ in terms of $x_k$:
      \begin{center}
        \begin{tikzpicture}
          \begin{axis}[xmin=0, xmax=22, ymin=0, ymax=20, x=0.06in,
            y=0.06in, hide y axis, axis x line=bottom, clip=false,
            xlabel=$x$, xtick=\empty, hide x axis] 
            \begin{scope}[yshift=8pt]
              \addplot+[domain=7.5:10, fill=cyan, draw=none,yshift=-8pt]
              {sqrt(20^2-\x^2)} -- (10, 0) -- (7.5, 0); %
              \draw (7.5, 0) -- (7.5, {sqrt(20^2-7.5^2)}); %
              \draw (10, 0) -- node[right, inner sep=1pt] {$h_k$} (10,
              {sqrt(20^2-10^2)}); %
              \draw (0, 0) -- (20, 0) arc[radius=20, start angle=0, end
              angle=90]; %
              \draw[blue, thick] (0, 0) -- (0, 20); %
              \draw[decorate, decoration={brace, mirror}, yshift=-2pt] (0,
              0) -- node[yshift=-2pt, below] {$x_k$} (10, 0); %
              \draw[red, thick] (0, 0) -- (10, {sqrt(20^2-10^2)}); %
            \end{scope}
          \end{axis}
        \end{tikzpicture}
      \end{center}
      We now see a right triangle in the picture: the hypotenuse is the
      radius of the pond, which is 100; the two legs have length $x_k$
      and $h_k$.  Therefore, by the Pythagorean Theorem, $x_k^2 + h_k^2 =
      20^2$, so $h_k = \sqrt{20^2 - x_k^2}$.  So, we can approximate the
      area of this slice as $h_k \Delta x = \sqrt{20^2 - x_k^2}\Delta x$.
      Therefore,
      \[ \text{number of fish in $k$-th slice} \approx \rho(x_k) \cdot \sqrt{20^2
        - x_k^2} \Delta x.\]
      Summing over all slices and taking the limit as the number of slices
      $\to \infty$, the number of fish in one quarter of the pond is
      $\displaystyle \int_0^{20} \rho(x) \cdot \sqrt{20^2 - x^2}\,dx$.
      Therefore, the total number of fish in the pond is
      \begin{align*}
        \boxed{4 \int_0^{20} \rho(x) \cdot \sqrt{20^2 - x^2}\,dx} &= 4
        \int_0^{100} (40 - x) \sqrt{20^2 - x^2}\,dx 
      \end{align*}
    \end{solution}

  \newpage 

 
\item
\begin{workshop}
    Students should have seen a very similar example of slicing a cone in class yesterday, but it nonetheless might be challenging for students. Highlights to include in the write-up for this question is how to approximate the volume of the $k$th slice (a.k.a. the $k$th disk) as the volume of a cylinder, and how to relate the height and radius of the cone using similar triangle geometry.
\end{workshop}
\begin{problem}
  A cone is 28 inches tall and has a radius of 4 inches at
  the top.  It is filled with layers of different types of
  fruit purees with varying densities.  The density of puree is given
  by $\rho(x)$ ounces per cubic inch, where $x$ is the height from the
  bottom (tip) of the cone. 
  
  Write a general Riemann sum that approximates the weight of fruit puree in the cone. Be sure to use proper summation notation and draw and label the dimensions of a typical slice. Then, write but do not evaluate a definite integral that gives the exact weight of fruit puree in the cone.
    \end{problem}

    \begin{solution}
    First, let's make sure we understand the meaning of $x$ in this
    problem.  Since $x$ represents distance from the bottom tip of the cone, $x = 0$ is at
    the bottom tip of the cone and $x = 28$ at the top.

    As always, the key in density problems is to slice so that the density
    within each slice is approximately constant.  This allows us to
    reasonably approximate the amount (of puree, in this case) in each
    slice by multiplying the density within the slice by the volume of the
    slice.  Since the density varies with height, we will slice the cone
    parallel to the top face of the cone.  This amounts to slicing the interval $[0, 8]$
    (on the $x$-axis).  As usual, we'll focus on the $k$-th slice (the
    whole slice is shown to the left):
    \begin{center}
      \begin{tikzpicture}[x=0.5cm, y=0.5cm]

        \begin{axis}[xmin=-4, xmax=4.5, ymin=0, ymax=8, clip=false, hide
          x axis, axis y line=right, x=0.5cm, y=0.5cm, ytick={0, 0.7, ...,
            7.01}, yticklabels={$x_0 = 0$,,,,,,$x_k$,,,,$x_n = 28$}]
            \coordinate (right) at (4, 7);
            \coordinate (left) at (-4, 7);
            \coordinate (tip) at (0, 0);
          \draw[fill=yellow, draw=none] (2, 3.5) arc[x radius=2,
          y radius=0.5, start angle=0, end angle=-180] -- (-2.4, 4.2)
          arc[x radius=2.4, y radius=0.6, start angle=180, end
          angle=360] -- cycle; %
          \draw (right) arc[x radius=4, y radius=1, start angle=0, end
          angle=180]; %
          \draw (right) -- (tip) -- (left); %
          \draw (4.5, 8) node[above] {$x$}; %
          \pgfplotsforeachungrouped \y in {0, 0.7, ..., 7.01} {
            \pgfmathsetmacro{\x}{4/7*\y} %
            \edef\temp{%
              \noexpand\draw (\x, \y) arc[x radius=\x, y radius={\x/4},
              start angle=0, end angle=-180];
            }              
            \temp
          }
          \draw [->,decorate, decoration={snake,amplitude=.4mm,segment
            length=2mm,post length=1mm}]  (-2.5, 3.75) -- (-5.5, 3.75);
          \begin{scope}[xshift=-4cm]
            \draw[fill=yellow, draw=none] (2, 3.5) arc[x radius=2, y
            radius=0.5, start angle=0, end angle=-180] -- (-2.4, 4.2)
            arc[x radius=2.4, y radius=0.6, start angle=180, end
            angle=360] -- cycle; %
            \draw[fill=yellow, draw=none] (2.4, 4.2) arc[x radius=2.4, y
            radius=0.6, start angle=0, end angle=360]; %
            \draw (2.4, 4.2) arc[x radius=2.4, y radius=0.6, start
            angle=0, end angle=360] -- (2, 3.5) arc[x radius=2, y
            radius=0.5, start angle=0, end angle=-180] -- (-2.4, 4.2); %
          \end{scope}
        \end{axis}
      \end{tikzpicture}
    \end{center}
    The $k$-th slice looks approximately like a disk of thickness $\Delta
    x$, so we can approximate its volume as $\pi (\text{radius})^2 \Delta
    x$.  However, we must express the radius in terms of $x_k$; we can do
    this using similar triangles.  Let's call the radius we want $r_k$: 
    \begin{center}
      \begin{tikzpicture}[x=0.5cm, y=0.5cm]
        

        \begin{axis}[xmin=-4, xmax=4.5, ymin=0, ymax=8, clip=false, hide
          x axis, axis y line=right, x=0.5cm, y=0.5cm, ytick={0, 0.7, ...,
            7.01}, yticklabels={$x_0 = 0$,,,,,,$x_k$,,,,$x_n = 8$}]
            \coordinate (right) at (4, 7);
        \coordinate (left) at (-4, 7);
        \coordinate (tip) at (0, 0);
          \draw[fill=yellow, draw=none] (2, 3.5) arc[x radius=2,
          y radius=0.5, start angle=0, end angle=-180] -- (-2.4, 4.2)
          arc[x radius=2.4, y radius=0.6, start angle=180, end
          angle=360] -- cycle; %
          \draw (right) arc[x radius=4, y radius=1, start angle=0, end
          angle=360]; %
          \draw (right) -- (tip) -- (left); %
          \draw (4.5, 8) node[above] {$x$}; %
          \draw[dashed] (2.4, 4.2) arc[x radius=2.4, y radius=0.6, start
          angle=0, end angle=180]; % 
          \pgfplotsforeachungrouped \y in {3.5, 4.2} {
            \pgfmathsetmacro{\x}{4/7*\y} %
            \edef\temp{%
              \noexpand\draw (\x, \y) arc[x radius=\x, y radius={\x/4},
              start angle=0, end angle=-180];
            }              
            \temp
          }
          \draw[red] (tip) -- ($0.5*(right)+0.5*(left)$) -- (right) --
          cycle;
          \draw[red] (0, 4.2) -- node[above, inner sep=0.5pt] {$r_k$}
          (2.4, 4.2); 
        \end{axis}
      \end{tikzpicture}
    \end{center}
    The larger triangle in the picture has height 28 and width 4; the
    smaller triangle has height $x_k$ and width $r_k$.  So, $\frac{4}{28} =
    \frac{r_k}{x_k}$, which means $r_k = \frac{4}{28} x_k = \frac{1}{7}
    x_k$.  Therefore,
    \begin{align*}
      \text{volume of $k$-th slice} &\approx \pi \left( \frac{1}{7} x_k
      \right)^2 \Delta x \\
      & = \frac{\pi}{49} x_k^2 \Delta x \text{ cubic inches} \\
      \text{density of fruit puree in $k$-th slice} & \approx \rho(x_k)
      \text{ ounces per cubic inch}\\
      \text{amount of fruit puree in $k$-th slice} & \approx \rho(x_k) \cdot 
      \frac{\pi}{49} x_k^2 \Delta x \\
      & = \frac{\pi}{49} x_k^2 \rho(x_k)
      \Delta x \text{ ounces}
    \end{align*}
    Summing over all slices, the amount of puree in the cup is
    approximately $\displaystyle \sum_{k=1}^n \frac{\pi}{49} x_k^2
    \rho(x_k) \Delta x$.  Taking the limit as $n \to \infty$, the amount
    of puree in the cone is exactly \fbox{$\displaystyle \int_0^{28} \frac{
        \pi}{49} x^2 \rho(x)\,dx$}.
  
\end{solution}

  
\pspace{\newpage}
\begin{workshop}
    The purpose of this question is to solidify the idea that a solution to a differential equation is a function. By the end of our second class on differential equations, students should be well acquainted with the differential equations listed in \#4c and \#4d; questions \#4a and \#4b are meant to remind students that we already know how to solve differential equations that are written exclusively in terms of $x$ using our integration skills!
\end{workshop}
\item
Find the general solution to each differential equation. Report the solutions in the form $y=f(x)$.
\begin{framed}
    A solution to a differential equation is a \underline{function} that ``solves" the differential equation, or makes the differential equation true. The \underline{general} solution to a differential equation involves an arbitrary constant $C$ that captures the family of functions that are solutions to the differential equation. 
\end{framed}
\begin{enumerate} 
\begin{multicols}{2}
    \item $ \dfrac{dy}{dx}=x$
    
    \begin{solution}
        $y=\dfrac{1}{2}x^2+C$
    \end{solution}
      \pspace{0.5in}
    \item $ \dfrac{dy}{dx}=kx$
    
    \begin{solution}
        $y=\dfrac{k}{2}x^2+C$
    \end{solution}
   \pspace{0.5in}
    \item $ \dfrac{dy}{dx}=y$

    \begin{solution}
        $y=Ce^x$
    \end{solution}
   \pspace{0.5in}
    \item $ \dfrac{dy}{dx}=ky$

    \begin{solution}
        $y=Ce^{kx}$
    \end{solution}
   \pspace{0.5in}
    \end{multicols}
      \pspace{0.5in}
\end{enumerate}
\begin{workshop}
Students completed a very similar problem in class, and this is meant as a review/revisit problem that acquaints them with the differential equation type $\dfrac{dy}{dx}=ky$.
\end{workshop}
\item  

 The rate at which a population grows is proportional to its size. 
  \begin{enumerate}
  \item
    \begin{problem}
      If $P(t)$ is the population size at time $t$, write a differential
      equation that describes the situation.      
    \end{problem}
\pspace{\vfill}
  \begin{framed}
     {$A$ is proportional to $B$} translates mathematically to $A = kB$ for $k$ a constant.
  \end{framed}
    \begin{solution}
      We can translate the information we're given into an equation: 
      \[ \underbrace{\text{The rate at which the population
          grows}}_{\frac{dP}{dt}} \text{ } \underbrace{\text{is
          proportional to}}_{{} = k \text{ times}} \text{ }
      \underbrace{\text{the population size}}_{P}.\] That is,
      $\boxed{\frac{dP}{dt} = k P}$. 
    \end{solution}

  \item
    \begin{problem}
      When the population is 800, it grows at a rate of 100 individuals per hour. Incorporate this into
      your differential equation.
    \end{problem}

    \begin{solution}
      We are told that, when $P = 800$, $\frac{dP}{dt} = 100$.  Plugging
      this into our differential equation $\frac{dP}{dt} = kP$, $100 = 800
      k$, so $k = \dfrac{1}{8}$.  Therefore, our differential equation becomes
      $\boxed{\frac{dP}{dt} = \dfrac{1}{8} P}$.
    \end{solution}
    \pspace{\vfill}
\item
    \begin{problem}
      Does it take more than an hour, less than an hour, or exactly one
      hour for the population to grow from 800 to 900?  Explain
      how you know.
    \end{problem}

    \begin{solution}
      We were told that, when the colony has a population of 800 bacteria,
      the colony is growing at an instantaneous rate of 100 bacteria per
      hour.  As the colony grows, its rate of growth also grows (because
      its rate of growth is proportional to the size of the colony).
      Thus, the colony will immediately start growing faster than 100
      bacteria per hour, so it will take \fbox{less than one hour} for the
      population to grow to 900.
    \end{solution}
 
 \pspace{\vfill}

\item
    \begin{problem}
    
    If the initial population is 400, find a
      formula for $P(t)$. 
    \end{problem}

    \begin{solution}
      The general solution to the differential equation $\frac{dP}{dt} =
      \frac{1}{8} P$ is $P(t) = C e^{\frac{1}{8} t}$.  Knowing the initial population is
      400 says that $P(0) = 400$, so $C e^{\frac{1}{8} (0)} = 400$, which tells us
      that $C = 400$.  So, $P(t) = \boxed{400 e^{\frac{1}{8} t}}$.
    \end{solution}
    \pspace{\vfill}
  \item
    \begin{problem}
      How long does it take the population to double in size?
    \end{problem}

    \begin{solution}
      Since the population starts out with 400 bacteria, we're looking for
      the time at which the population reaches 800 bacteria.  Since the
      population at time $t$ is $400 e^{\frac{1}{8} t}$, we solve
      \begin{align*}
        400 e^{0.2 t} &= 800 \\
        e^{\frac{1}{8} t} &= 2 \\
        \frac{1}{8} t &= \ln 2 \\
        t &= 8 \ln 2
      \end{align*}
      So, it takes $\boxed{8 \ln 2}$ hours for the population to double
      (this is about $5.55$ hours). 
    \end{solution}
    \pspace{\vfill}
  \end{enumerate}
  \pspace{\newpage}
  \begin{workshop}
      The purpose of this question is to remind students that, while we have not yet learned the methods to solve most types of differential equations, we can check a proposed solution to differential equations, in part, by differentiating that solution.
  \end{workshop}
  \item 
  \begin{enumerate}
      \item Verify that $y=2e^{-x}+x-1$ is a solution to the differential equation $\frac{dy}{dx}=x-y$. 
      
      \begin{solution}
          Given the formula for $y$, 
          \begin{align*}
              \frac{dy}{dx}&=\fbox{$-2e^{-x}+1$}.
          \end{align*}
          We also know that $\frac{dy}{dx}=x-y$. Plugging in $y$, 
          \begin{align*}
            \frac{dy}{dx}&=x-(2e^{-x}+x-1)\\&=x-2e^{-x}-x+1\\&=\fbox{$-2e^{-x}+1$}
          \end{align*}.
      \end{solution}
      \pspace{\vfill}
      \item Verify that $y=\dfrac{1}{1-x}$ is a solution to the differential equation $\frac{dy}{dx}=y^2$. 
      
      \begin{solution}
          Let's rewrite $y=\dfrac{1}{1-x}$ as $y=(1-x)^{-1}$.
          Given this formula for $y$, 
          \begin{align*}
              \frac{dy}{dx}&=\left[-(1-x)^{-2}\right]\cdot(-1)\\&=(1-x)^{-2}\\&=\dfrac{1}{(1-x)^2}\\&=\fbox{$\left(\dfrac{1}{1-x}\right)^2$}.
          \end{align*}
          We also know that $\frac{dy}{dx}=y^2$. Plugging in $y$, 
          \begin{align*}
            \frac{dy}{dx}&=\fbox{$\left(\dfrac{1}{1-x}\right)^2$}
          \end{align*}.
      \end{solution}
      \pspace{\vfill}
      \item Verify that $y=e^{3x}-\dfrac{1}{2}e^x$ is a solution to the differential equation $\frac{dy}{dx}=3y+e^x$.

      \begin{solution}
      Given the formula for $y$,
      \begin{align*}
              \frac{dy}{dx}=\fbox{$3e^{3x}-\dfrac{1}{2}e^x$}.
          \end{align*}
          We also know that $\frac{dy}{dx}=3y+e^x$. Plugging in $y$, 
          \begin{align*}
            \frac{dy}{dx}&=3(e^{3x}-\dfrac{1}{2}e^x)+e^x \\ &= 3e^{3x}-\dfrac{3}{2}e^x+e^x\\&=\fbox{$3e^{3x}-\dfrac{1}{2}e^x$}
          \end{align*}.
      \end{solution}
      \pspace{\vfill}
      \item Verify that $y=e^x+\dfrac{1}{2}\sin x-\dfrac{1}{2}\cos x$ is a solution to the differential equation $\frac{dy}{dx}=\cos x + y$. 
      
      \begin{solution}
      Given the formula for $y$,
      \begin{align*}
              \frac{dy}{dx}=\fbox{$e^x+\dfrac{1}{2}\cos x+\dfrac{1}{2}\sin x$}.
          \end{align*}
          We also know that $\frac{dy}{dx}=\cos x + y$. Plugging in $y$, 
          \begin{align*}
            \frac{dy}{dx}&=\cos x + (e^x+\dfrac{1}{2}\sin x-\dfrac{1}{2}\cos x)  \\ &=\fbox{$e^x+\dfrac{1}{2}\sin x+\dfrac{1}{2}\cos x$}
          \end{align*}.
      \end{solution}
      \pspace{\vfill}
  \end{enumerate}
 \end{enumerate}
\end{document}